\documentclass{beamer}
% Thème moderne et sobre
\usetheme[progressbar=frametitle]{metropolis}
\usecolortheme{spruce}
\usepackage{tikz}
\usetikzlibrary{positioning}
\usepackage[utf8]{inputenc}
\usepackage[T1]{fontenc}
\usepackage[french]{babel}
\usepackage{graphicx}
\usepackage{hyperref}

% Pour encadrés visuels
\setbeamercolor{block title alerted}{fg=white,bg=orange!85!black}
\setbeamercolor{block body alerted}{fg=black,bg=orange!10!white}
\setbeamercolor{block title example}{fg=white,bg=teal!80!black}
\setbeamercolor{block body example}{fg=black,bg=teal!10!white}

\title[Conception d'une Ontologie et Moteur d'Inférence]{Présentation du Travail de Conception :\newline Modules et Moteur d'Inférence}
\author{Votre Nom}
\institute{Votre Institution}
\date{26 Mai 2026}

\begin{document}

% Slide d'ouverture visuelle
\begin{frame}[plain]
  \vfill
  \centering
  {\LARGE\textbf{« L’agriculture de demain se construit aujourd’hui par la connaissance et l’interopérabilité. »}}
  \vfill
\end{frame}

\begin{frame}
  \titlepage
\end{frame}

% Sommaire
\begin{frame}{Plan}
  \tableofcontents
  \vspace{0.5cm}
  \begin{center}
    \includegraphics[width=0.25\linewidth]{architecture.png}
  \end{center}
  \pause
  \alert{\small{Chaque module et scénario sera illustré en détail.}}
\end{frame}

\section{Introduction}
\begin{frame}{Pourquoi une ontologie agricole ?}
  \begin{block}{Enjeux}
    \begin{itemize}
      \item Complexité croissante des systèmes agricoles
      \item Besoin d’interopérabilité et de partage des connaissances
      \item Valorisation des données pour l’innovation
    \end{itemize}
  \end{block}
\end{frame}
\begin{frame}{Contexte et Objectifs}
  \begin{itemize}
    \item Besoin de modélisation des connaissances dans le domaine agricole
    \item Objectif : Concevoir une ontologie robuste et un moteur d'inférence performant
    \item Valorisation des données pour l'aide à la décision
  \end{itemize}
\end{frame}

\section{Méthodologie de Conception}
\begin{frame}{Méthodes de Conception}
  \begin{enumerate}
    \item Analyse documentaire et entretiens avec des experts
    \item Construction itérative de l’ontologie
    \item Validation croisée des modules et des scénarios
    \item Tests d’interopérabilité et d’intégration
  \end{enumerate}
  \pause
  \begin{alertblock}{Clé de réussite}
    L’implication continue des experts métiers.
  \end{alertblock}
\end{frame}
\begin{frame}{Démarche de Conception}
  \begin{enumerate}
    \item Analyse des besoins et collecte des exigences
    \item Modélisation des concepts clés (climat, sol, culture, pratiques)
    \item Définition des relations et propriétés
    \item Validation par des experts du domaine
  \end{enumerate}
\end{frame}

\begin{frame}{Outils et Technologies}
  \begin{itemize}
    \item OWL, RDF, Protégé
    \item Raisonneurs : Pellet, Hermit
    \item Intégration avec des modules Python/Java pour l'inférence
  \end{itemize}
\end{frame}

\section{Architecture des Modules}

\begin{frame}{Schéma conceptuel de l’architecture}
  \centering
  \begin{tikzpicture}[node distance=2cm, every node/.style={draw, fill=teal!10, rounded corners, minimum width=2.5cm, minimum height=1cm, align=center}]
    \node (climat) {Climat};
    \node (sol) [right=of climat] {Sol};
    \node (culture) [below=of climat] {Culture};
    \node (pratique) [right=of culture] {Pratique};
    \node (situation) [below=1.5cm of culture, fill=orange!20] {\textbf{Situation (Interopérabilité)}};
    \draw[->, thick] (climat) -- (situation);
    \draw[->, thick] (sol) -- (situation);
    \draw[->, thick] (culture) -- (situation);
    \draw[->, thick] (pratique) -- (situation);
  \end{tikzpicture}
\end{frame}

\begin{frame}{Vue d'Ensemble des Modules}
  \begin{figure}[h]
    \centering
    \includegraphics[width=0.8\linewidth]{architecture.png}
    \caption{Architecture modulaire du système}
  \end{figure}
\end{frame}

\begin{frame}{Module Climat}
  \begin{itemize}
    \item Gestion des données météorologiques (température, précipitations, humidité)
    \item Intégration de sources climatiques variées
    \item Impact sur les recommandations agronomiques
  \end{itemize}
\end{frame}

\begin{frame}{Module Sol}
  \begin{itemize}
    \item Caractérisation des types de sols (texture, pH, fertilité)
    \item Prise en compte des propriétés physiques et chimiques
    \item Influence sur le choix des cultures et pratiques
  \end{itemize}
\end{frame}

\begin{frame}{Module Culture (Crop)}
  \begin{itemize}
    \item Modélisation des espèces cultivées
    \item Cycle de vie, besoins spécifiques
    \item Adaptation aux conditions locales
  \end{itemize}
\end{frame}

\begin{frame}{Module Pratique}
  \begin{itemize}
    \item Représentation des pratiques agricoles (rotation, fertilisation, irrigation)
    \item Effets sur la durabilité et la productivité
  \end{itemize}
\end{frame}

\begin{frame}{Module d'Interopérabilité : Situation (situation.rdf)}
  \begin{itemize}
    \item Centralise l'information sur les situations agronomiques
    \item Permet l'interconnexion entre les modules (sol, climat, culture, pratique)
    \item Facilite l'intégration de nouveaux scénarios et l'interopérabilité des données
    \item Base pour l'inférence et la personnalisation des recommandations
  \end{itemize}
  \pause
  \begin{exampleblock}{Focus}
    Le module \textbf{Situation} est le cœur de l’architecture : il permet de croiser toutes les dimensions pour générer des recommandations adaptées à chaque contexte agronomique.
  \end{exampleblock}
\end{frame}


% Présentation détaillée des scénarios de situations agronomiques
\section{Scénarios de Situations Agronomiques}

\begin{frame}{Résumé graphique des scénarios}
  \centering
  \begin{tabular}{|c|c|c|c|c|}
    \hline
    \textbf{\#} & \textbf{Sol} & \textbf{Climat} & \textbf{Culture} & \textbf{Pratique} \\
    \hline
    1 & Argileux & Humide & Riz & Rotation annuelle \\
    2 & Sableux & Sec & Mil & Irrigation min. \\
    3 & Limoneux & Tempéré & Blé & Fertilisation org. \\
    4 & Acide & Tropical & Manioc & Jachère \\
    5 & Calcaire & Méditerranéen & Vigne & Taille hivernale \\
    \hline
    ... & ... & ... & ... & ... \\
    \hline
  \end{tabular}
  \vspace{0.3cm}
  \alert{\small{Tous les scénarios sont modélisés dans situation.rdf}}
\end{frame}

\begin{frame}{Méthodes de Conception}
  \begin{enumerate}
    \item Analyse documentaire et entretiens avec des experts
    \item Construction itérative de l'ontologie
    \item Validation croisée des modules et des scénarios
    \item Tests d'interopérabilité et d'intégration
  \end{enumerate}
  \pause
  \begin{block}{Points forts}
    \begin{itemize}
      \item Approche collaborative et évolutive
      \item Validation continue par des cas réels
    \end{itemize}
  \end{block}
\end{frame}

\section{Moteur d'Inférence}
\begin{frame}{Fonctionnement du Moteur d'Inférence}
  \begin{itemize}
    \item Règles basées sur l'ontologie
    \item Raisonnement automatique pour l'aide à la décision
    \item Exemples de règles :
    \begin{itemize}
      \item Si le sol est argileux et le climat humide, alors recommander la culture du riz
      \item Si la pratique est la rotation, alors améliorer la fertilité du sol
    \end{itemize}
  \end{itemize}
  \pause
  \begin{alertblock}{Atout}
    Le moteur d’inférence permet d’automatiser la personnalisation des recommandations pour chaque situation.
  \end{alertblock}
\end{frame}

\begin{frame}{Exemple de Chaîne d'Inférence}
  \begin{block}{Cas d'utilisation}
    \begin{itemize}
      \item Entrée : Données sur le sol, le climat et la culture
      \item Traitement : Application des règles d'inférence
      \item Sortie : Recommandations personnalisées
    \end{itemize}
  \end{block}
  \pause
  \begin{exampleblock}{Visualisation}
    \centering
    %\includegraphics[width=0.5\linewidth]{inference_chain.png}
    %\newline
    \small{Chaîne d’inférence illustrant le passage des données à la recommandation. (Image à insérer ici)}
  \end{exampleblock}
\end{frame}

\section{Résultats et Validation}

\begin{frame}{Résultats Obtenus}
  \begin{block}{Validation}
    \begin{itemize}
      \item Ontologie validée par des experts du domaine
      \item Interopérabilité assurée entre les modules via situation.rdf
      \item Génération automatique de recommandations adaptées à chaque scénario
    \end{itemize}
  \end{block}
  \pause
  \begin{alertblock}{Impact}
    \begin{itemize}
      \item Amélioration de la prise de décision agronomique
      \item Facilité d’extension et d’intégration de nouveaux modules/scénarios
      \item Base solide pour des applications d’aide à la décision
    \end{itemize}
  \end{alertblock}
\end{frame}

\section{Conclusion}

\begin{frame}{Bilan et Perspectives}
  \begin{block}{Bilan}
    \begin{itemize}
      \item Ontologie modulaire et évolutive
      \item Moteur d’inférence adaptable
    \end{itemize}
  \end{block}
  \pause
  \begin{exampleblock}{Perspectives}
    \begin{itemize}
      \item Extension à d’autres domaines agricoles
      \item Intégration de données massives (big data, IoT)
      \item Déploiement sur plateformes d’aide à la décision
    \end{itemize}
  \end{exampleblock}
\end{frame}

\begin{frame}{Questions ?}
  \centering
  \Huge{Merci pour votre attention !}
\end{frame}

\end{document}
